  \begin{block}{Motivation: Trackers on Essential Public Services}

    Citizens rely on government websites for essential services such as applying for a passport, filing taxes, or obtaining healthcare information, and they generally expect privacy. Unlike commercial platforms, government portals are often the \textbf{only available option}, leaving individuals with no ability to opt out. In practice, many government websites embed third-party trackers that collect browsing behavior, unique identifiers, and digital fingerprints, raising privacy, security, and data sovereignty concerns, including \textbf{cross-border tracker-data flows}. Prior work mostly targets commercial websites or small, single-country government samples and is rarely longitudinal. We present the \textbf{first large-scale, global, longitudinal study} of tracking on government websites.

  \end{block}
